\section{A reversible design and finite horizons}\label{sec:reversible}

This section is not part of the journal version. It reads a reversible generator off the
first-order modal-support map by a design fixed in advance
(Theorem~\ref{thm:reversible-reconstruction}), counts the dimension
(Proposition~\ref{thm:minimality}), bounds the error (Theorem~\ref{thm:stability}) and treats a
finite horizon uniformly over a compact class (Theorem~\ref{thm:finite-horizon}).

\subsection{The design}\label{sec:reversible-reconstruction}

For a reversible chain the modal supports determine all of $Q$, and an explicit finite
experiment reads it. Let $Q$ be reversible and irreducible with stationary law $\pi$, on a known
connected graph $(V,E)$ with $m=|E|$ edges, where $\{i,j\}\in E$ exactly when $q_{ij}>0$. We
call it the graph $E$.
Put $w_{ij}=q_{ij}q_{ji}$ for the edge products and
\[
 S=D_{\sqrt\pi}(-Q)D_{\sqrt\pi}^{-1},\qquad S_{ii}=q_i,\qquad S_{ij}=-\sqrt{w_{ij}}\ \ (\{i,j\}\in E),
\]
the symmetrized matrix. We assume known constants $R,b>0$ with $\|S\|_2\le R$ and
$\sqrt{w_{ij}}\ge b$ on every edge. The field and the time are calibrated.

We use 2 facts throughout. First, $S$ is symmetric and similar to $-Q$, whose eigenvalues
have nonnegative real parts, so $S$ is positive semidefinite. Hence $0\le q_k\le R$ for every
$k$, $\sqrt{w_{kl}}\le R$ on every edge, and $|q_k-q_l|\le R$. Second, for a face $F$ on which
$E$ induces a connected graph, Proposition~\ref{prop:variational}(c) gives
\begin{equation}\label{eq:lambda-sym}
 \lambda_F(v)=\lambda_{\min}\bigl(S_F-\operatorname{diag}v_F\bigr).
\end{equation}
Since $Q$ has symmetric support, the admissible faces are exactly these connected faces
(Lemma~\ref{lem:admissible}).

The design uses, for a retained pair $(i,j)$, the field
\[
 v_i=\xi,\qquad v_j=0,\qquad v_k=-\bar v\ \ (k\ne i,j),
\]
which suppresses every other state, with the constants
\[
 \bar\xi=2R,\qquad \tau_0=\frac1{2R},\qquad \tau_{\max}=\frac{8R}{b^2},\qquad \bar v=4R+4\tau_{\max}R^2,
\]
and the ranges $|\xi|\le3R$ and $\frac1{4R}\le\tau\le \tau_{\max}$. Every constant is fixed by $R$ and $b$.

\begin{lemma}[suppression in the reversible design]\label{lem:suppression}
Under the design for the pair $(i,j)$, every irreducible face that contains a suppressed state
has score at least $\frac34$ larger than the best face inside $\{i,j\}$, uniformly on the ranges
of $\xi$ and $\tau$. Hence the strict winner is one of $\{i\}$, $\{j\}$ and $\{i,j\}$, and the last
only if $\{i,j\}$ is an edge.
\end{lemma}

\begin{proof}
Let $F$ be connected and let $A=F\setminus\{i,j\}$ be nonempty.

\emph{Wholly suppressed faces.} If $F\cap\{i,j\}=\emptyset$, then
$S_F-\operatorname{diag}v_F=S_F+\bar vI\ge\bar vI$, so $\lambda_F(v)\ge\bar v$
by~\eqref{eq:lambda-sym}. The singleton $\{j\}$ has score $\tau q_j\le\tau R$. Hence
$\Phi_{\tau,v}(F)-\Phi_{\tau,v}(\{j\})\ge\tau(\bar v-R)\ge(3R+4\tau_{\max}R^2)/(4R)\ge\frac34$.

\emph{Mixed faces.} Let $G_0=F\cap\{i,j\}$ be nonempty. In the order $G_0$, $A$,
\[
 S_F-\operatorname{diag}v_F=\begin{pmatrix}B&Z\\Z^\top&S_A+\bar vI\end{pmatrix},\qquad
 B=S_{G_0}-\operatorname{diag}v_{G_0},\quad Z=S[G_0,A].
\]
Here $\|Z\|_2\le R$ and $S_A\ge0$. Let $e=\lambda_{\min}(B)$, so
$|e|\le\|S_{G_0}\|_2+|\xi|\le4R$. The test vector $(x,0)$, with $x$ a unit eigenvector of $B$
for $e$, gives $\lambda_F(v)\le e$. For a unit vector $(x,y)$,
\[
 \Bigl\langle\binom xy,(S_F-\operatorname{diag}v_F)\binom xy\Bigr\rangle
 \ge e|x|^2-2R|x||y|+\bar v|y|^2\ge\lambda_{\min}\begin{pmatrix}e&-R\\-R&\bar v\end{pmatrix}
 \ge e-\frac{R^2}{\bar v-e},
\]
by $\sqrt{u^2+R^2}\le u+R^2/(2u)$ with $u=(\bar v-e)/2>0$. Since $\bar v-e\ge\bar v-4R=4\tau_{\max}R^2$,
we get $e-1/(4\tau_{\max})\le\lambda_F(v)\le e$. Now $e=\lambda_{G_*}(v)$ for an irreducible face
$G_*\subseteq\{i,j\}$. If $G_0$ is connected, take $G_*=G_0$ and use~\eqref{eq:lambda-sym}. If
$G_0=\{i,j\}$ is not an edge, $B$ is diagonal and $G_*$ is the singleton with the smaller
diagonal entry. In both cases $|G_*|\le|F|-1$, so
$\Phi_{\tau,v}(F)-\Phi_{\tau,v}(G_*)\ge1-\tau/(4\tau_{\max})\ge\frac34$ as $\tau\le \tau_{\max}$.

So the unique minimizer of $\Phi_{\tau,v}$, if there is one, is an irreducible face inside
$\{i,j\}$, and Proposition~\ref{prop:controlled-score} identifies it with the strict winner.
\end{proof}

\begin{theorem}[reversible reconstruction]\label{thm:reversible-reconstruction}
Let $Q$ be reversible and irreducible on the known connected graph $E$ with $\|S\|_2\le R$ and
$\sqrt{w_{ij}}\ge b$ on edges, and let the start have full support.
\begin{enumerate}[(a)]
\item (Exit differences.) For each $i\ne1$, under the design for the pair $(i,1)$ at
$\tau=\tau_0$, the strict winner changes from $\{1\}$ to $\{i\}$ at $\xi=x_i=q_i-q_1$, and
$|x_i|\le R$. Put $x_1=0$.
\item (Products.) For each edge $\{i,j\}$, under the design for $(i,j)$ at $\xi=\bar\xi$, the
strict winner changes from a singleton to the pair at
\[
 \tau=t_{ij}=\frac{\sqrt{s_{ij}^2+4w_{ij}}+s_{ij}}{2w_{ij}},\qquad
 s_{ij}=\bar\xi-x_i+x_j,
\]
with $1/R\le t_{ij}<\tau_{\max}/2$, and $w_{ij}=s_{ij}/t_{ij}+t_{ij}^{-2}$.
\item (Reconstruction.) Let $B$ be the symmetric matrix with diagonal $x_i$, entries
$-\sqrt{w_{ij}}$ on the edges and zero elsewhere, let $\lambda=\lambda_{\min}(B)$ and let $\phi$ be
its positive unit eigenvector. Then
\[
 S=B-\lambda I,\qquad Q=-D_\phi^{-1}SD_\phi,\qquad \pi_i=\phi_i^2 .
\]
So the $m+d-1$ limiting boundary coordinates $(x_i)_{i\ne1}$ and $(t_{ij})_{\{i,j\}\in E}$
determine $Q$, and the design is fixed by $R$, $b$ and the labels before any boundary is
observed.
\end{enumerate}
\end{theorem}

\begin{proof}
By Lemma~\ref{lem:suppression} and Proposition~\ref{prop:controlled-score}, at each point of the
ranges the strict winner is the unique minimizer of $\Phi_{\tau,v}$ among the irreducible faces
inside the retained pair, when it is unique. Here $\tau_0\in[\frac1{4R},\tau_{\max}]$, since $b\le R$.

(a) The candidates are $\{i\}$, $\{1\}$ and, if $\{i,1\}$ is an edge, the pair. By Weyl's
inequality $\lambda_{\{i,1\}}(v)\ge\min(q_i-\xi,q_1)-\sqrt{w_{i1}}$, so at $\tau_0$ the score of
the pair exceeds the better singleton score by at least $1-\tau_0R=\frac12$, and the pair never
wins. The singleton $\{i\}$ has the smaller score exactly when $q_i-\xi<q_1$, that is when
$\xi>x_i$. Since $|x_i|\le R$, both sides of $x_i$ lie in $[-3R,3R]$. Hence $\{1\}$ wins strictly
on $[-3R,x_i)$, $\{i\}$ wins strictly on $(x_i,3R]$, and the boundary is $\xi=x_i$.

(b) Put $p=q_i-\bar\xi$ and $s=s_{ij}$. Since $x_i-x_j=q_i-q_j$ by (a), $s=q_j-p$. Then
$R\le s\le3R$, and $\{i\}$ beats $\{j\}$ by $\tau s>0$. By~\eqref{eq:lambda-sym} the
pair has $\lambda_{\{i,j\}}(v)=p-g$, the smaller eigenvalue of
$\bigl(\begin{smallmatrix}p&-\sqrt w\\-\sqrt w&q_j\end{smallmatrix}\bigr)$ with $w=w_{ij}$,
where
\[
 g=\frac{\sqrt{s^2+4w}-s}2=\frac{2w}{\sqrt{s^2+4w}+s}\in(0,\sqrt w\,] .
\]
So $\Phi(\{i,j\})-\Phi(\{i\})=1-\tau g$. Hence $\{i\}$ wins strictly below $t_{ij}=1/g$, the pair
wins strictly above it, and $1/g$ is the displayed formula. Also $g\le\sqrt w\le R$ gives
$t_{ij}\ge1/R>\frac1{4R}$. The formula increases in $s$ and decreases in $w$, and $s\le3R$,
$w\ge b^2$, $b\le R$, so $t_{ij}\le(\sqrt{13}+3)R/(2b^2)<\tau_{\max}/2$. At $\tau=t_{ij}$ the pair
eigenvalue $p-1/t_{ij}$ is a root of $(p-y)(q_j-y)=w$, which gives $t_{ij}^{-1}(s+t_{ij}^{-1})=w$,
the inversion formula.

(c) $B=S-q_1I$, since $S$ has diagonal $q_k$, entries $-\sqrt{w_{kl}}$ on edges and zero off the
graph. The positive unit vector $\sqrt\pi$ satisfies $S\sqrt\pi=D_{\sqrt\pi}(-Q)\ones=0$. As
$S\ge0$, the eigenvalue $0$ is $\lambda_{\min}(S)$, and it is simple with eigenvector $\sqrt\pi$
by the theorem of Perron and Frobenius for the irreducible nonnegative matrix $RI-S$. So
$\lambda=-q_1$, $S=B-\lambda I$, $\phi=\sqrt\pi$ and $\pi_i=\phi_i^2$, and
$Q=-D_\phi^{-1}SD_\phi$ is the definition of $S$. The coordinates used are $d-1$ values $x_i$
and $m$ values $t_{ij}$, and the fields and time scales of the design depend only on $R$, $b$
and the labels.
\end{proof}

The reconstruction of a reversible generator from its exit rates and edge products is
classical. Kim and Landim~\cite[Lemma~5.1 and Corollary~5.3]{kl2024} recover the same
quantities from the values of $J_V$, and with the stationary law from $J_V^{-1}(0)$ (their
Lemma~5.2) deduce that $J_V$ determines a reversible generator. Here the coordinates are read
off the boundaries of strict regions of the weighted composition law, with a design fixed in
advance.

\begin{proposition}[the dimension count]\label{thm:minimality}
The reversible generators on the fixed connected graph $E$ form a real-analytic manifold of
dimension $m+d-1$, with global coordinates the edge products $(w_{ij})$ and the exit differences
$(q_i-q_1)$. Let $W$ be a nonempty open subset of it, for example the set where $\|S\|_2<R$ and
$\sqrt{w_{ij}}>b$ when that set is nonempty. Let $k<m+d-1$, let $f\colon W\to\R^k$ be $C^1$ and
let $\ell\colon W\to\Lambda$ be locally constant with $\Lambda$ finite. Then the observation
$(f,\ell)$ is not injective on any nonempty open subset of $W$. So no family of fewer than
$m+d-1$ locally $C^1$ scalar observations with locally constant finite labels is locally
identifying.
\end{proposition}

\begin{proof}
Let $\mathcal Q_E$ be the set of reversible generators with $q_{ij}>0$ exactly on the edges in
$E$, and $\kappa(Q)=((w_{ij})_E,(q_i-q_1)_{i\ne1})$. Given $w\in(0,\infty)^m$ and
$x\in\R^{d-1}$ (with $x_1=0$), form $B$ as in Theorem~\ref{thm:reversible-reconstruction}(c),
let $\lambda$ and $\phi$ be its least eigenvalue and positive unit eigenvector, which is simple
since $-B$ plus a large multiple of $I$ is irreducible and nonnegative, and put
$\Psi(w,x)=-D_\phi^{-1}(B-\lambda I)D_\phi$. This is a reversible generator on $E$, with respect
to $\phi^2$, with products $w_{ij}$ and exit rates $x_i-\lambda$, so $\kappa\circ\Psi=\mathrm{id}$,
and Theorem~\ref{thm:reversible-reconstruction}(c), whose proof uses only reversibility and
irreducibility, gives $\Psi\circ\kappa=\mathrm{id}$. The map $\Psi$ is real-analytic by the
implicit function theorem at a simple eigenvalue, and $\kappa$ is polynomial. So $\kappa$ is a
global real-analytic chart of $\mathcal Q_E$, of dimension $m+d-1$.

For the lower bound let $W'\subseteq W$ be nonempty and open. Shrinking it, we may assume that
$\ell$ is constant on $W'$ and that $W'$ is a ball in the chart $\kappa$, so $f$ is a $C^1$ map on
a ball of $\R^p$, $p=m+d-1$. Let $\varrho_f$ be the largest rank of $df$ on $W'$, attained at
$y_0$. The rank is lower semicontinuous, so it equals $\varrho_f$ near $y_0$, and by the constant
rank theorem $f$ is constant on a $C^1$ submanifold through $y_0$ of dimension
$p-\varrho_f\ge p-k\ge1$. So $(f,\ell)$ is not injective on $W'$.
\end{proof}

The chart is the continuous-time version of that of Wolfer and
Watanabe~\cite[Theorems~3, 5 and~6]{wolferwatanabe2021} for reversible kernels, and the lower
bound is the classical rank argument. Regularity is needed, since interleaving the digits of
the $m+d-1$ coordinates gives one injective scalar.

\begin{remark}[partial data]\label{rem:partial-data}
Suppose that the products are known on $J\subseteq E$ and the exit differences along the edges
of a comparison graph $(V,E')$. In the chart $\kappa$ the generators with these data have the
given $w_e$ on $J$ and any $w_e>0$ on $E\setminus J$. Their $x$ is the given one plus a vector
that is constant on each component of $(V,E')$ and zero on that of the state $1$. So they
are parametrized by $(0,\infty)^{m-|J|}\times\R^{\operatorname{comp}(E')-1}$, with
$\operatorname{comp}(E')$ the number of components of $(V,E')$, and the data identify $Q$ exactly
when $J=E$ and $(V,E')$ is connected. For instance, for
$\ell>\sqrt2$ and $z=\ell-2/\ell$ the generators
\[
 Q_\ell=\begin{pmatrix}-\ell&2/\ell&z\\ \ell/2&-\ell&\ell/2\\ z&2/\ell&-\ell\end{pmatrix}
\]
satisfy detailed balance with $\nu=(1,4/\ell^2,1)$, have all exit rates equal to $\ell$ and
$w_{12}=w_{23}=1$, and have $w_{13}=z^2$. So the data of a spanning tree miss a parameter.
\end{remark}

\begin{theorem}[stability]\label{thm:stability}
In the setting of Theorem~\ref{thm:reversible-reconstruction}, let the observed dimensionless
coordinates $x_i/R$ and $Rt_{ij}$ have sup-norm errors at most $\varepsilon$, and let
$\widehat Q$ be the reconstruction from the observed values. There are $\varepsilon_0$ and $C$
depending only on $d$ and $b/R$ such that for $\varepsilon\le\varepsilon_0$ the reconstruction
$\widehat Q$ is a reversible generator on $E$ and
\[
 \|\widehat Q-Q\|_2\le C\,R\,\varepsilon .
\]
Among the constants are the bound $\phi_{\min}\ge \phi_*=(b/R)^{d-1}/\sqrt d$ for the ground
vector $\phi=\sqrt\pi$ and the bound $\lambda_2(S)\ge g_*=2b\phi_*^2/(d-1)$ for the gap of $S$.
\end{theorem}

\begin{proof}
The reconstruction commutes with the scaling $Q\mapsto Q/R$, so we take $R=1$, $\bar\xi=2$ and
$b\le1$. Then $\sigma_{ij}=t_{ij}\ge1$ and $s_{ij}=2-x_i+x_j\in[1,3]$, by
Theorem~\ref{thm:reversible-reconstruction}. Hats denote reconstructed values, and
$\widehat x_1=x_1=0$.

\emph{Products.} Let $\varepsilon\le b^2/20$. Then $\widehat\sigma\ge\frac34$ and
$|\widehat\sigma^{-1}-\sigma^{-1}|\le\frac43\varepsilon$, and $w=s/\sigma+\sigma^{-2}$ gives
$|\widehat w-w|\le(\frac{20}3+\frac{28}9)\varepsilon<10\varepsilon$. So $\widehat w\ge b^2/2$ and
$|\sqrt{\widehat w}-\sqrt w|\le10\varepsilon/b$, and the largest absolute row sum gives
$\|\widehat B-B\|_2\le\beta:=10d\varepsilon/b$. By Weyl's inequality
$|\widehat\lambda-\lambda|\le\beta$, so $\|\widehat S-S\|_2\le2\beta$. The matrix $\widehat B$
has negative entries exactly on the edges in $E$, so the chart $\Psi$ of the proof of
Proposition~\ref{thm:minimality} makes $\widehat Q$ a reversible generator with the graph $E$.

\emph{The constants.} Row $l$ of $S\phi=0$ reads $q_l\phi_l=\sum_{k\sim l}\sqrt{w_{lk}}\phi_k$, so
$\phi_l\ge b\phi_k$ on every edge. Some $\phi_{k_0}$ is at least $1/\sqrt d$, and every state is
joined to $k_0$ by at most $d-1$ edges, so $\phi_{\min}\ge\phi_*$. For $y\perp\phi$ write
$y=\phi\circ f$. Then $\langle y,Sy\rangle=\sum_{\{k,l\}\in E}\chi_{kl}(f_k-f_l)^2$ with
$\chi_{kl}=\sqrt{w_{kl}}\phi_k\phi_l\ge b\phi_*^2$, while $\|y\|^2=\sum_k\pi_kf_k^2$ and
$\sum_k\pi_kf_k=0$. In a spanning tree $\mathcal T$ of $E$ the path from $k$ to $l$ has at most
$d-1$ edges, so the inequality of Cauchy and Schwarz gives
\[
 \|y\|^2=\tfrac12\sum_{k,l}\pi_k\pi_l(f_k-f_l)^2\le\tfrac{d-1}2\sum_{e\in\mathcal T}(\nabla_ef)^2
 \le\frac{d-1}{2b\phi_*^2}\langle y,Sy\rangle,
\]
and $\lambda_2(S)\ge g_*$.

\emph{The generator.} Write $\widehat\phi=a\phi+z$ with $z\perp\phi$ and $a>0$. Since
$Sz=(S-\widehat S)\widehat\phi$, we get $\|z\|\le2\beta/g_*$, and since $1-a\le1-a^2=\|z\|^2$,
also $\gamma=\|\widehat\phi-\phi\|\le\sqrt2\|z\|\le2\sqrt2\beta/g_*$. This is the bound of Davis
and Kahan in its simplest form. Let $\varepsilon$ be so small that $\gamma\le\phi_*/2$. Then
$\widehat\phi_{\min}\ge\phi_*/2$,
$|\widehat\phi_l/\widehat\phi_k-\phi_l/\phi_k|\le4\gamma/\phi_*^2$ and
$\widehat\phi_l/\widehat\phi_k\le2/\phi_*$, so on an edge
\[
 |\widehat q_{kl}-q_{kl}|
 =\bigl|\sqrt{\widehat w_{kl}}\,\widehat\phi_l/\widehat\phi_k-\sqrt{w_{kl}}\,\phi_l/\phi_k\bigr|
 \le\frac{20\varepsilon}{b\phi_*}+\frac{4\gamma}{\phi_*^2}.
\]
The diagonal satisfies $|\widehat q_k-q_k|=|\widehat S_{kk}-S_{kk}|\le2\beta$, and off the graph
both matrices vanish. So every entry of $\widehat Q-Q$ is at most $C_1(d,b)\varepsilon$, and
$\|\widehat Q-Q\|_2\le dC_1\varepsilon$. The conditions on $\varepsilon$ depend only on $d$ and
$b$.
\end{proof}

The perturbation bounds are standard. The point is that the constants come from $R$ and $b$
alone, with no separate assumption on the spectral gap.

\subsection{Finite horizons}\label{sec:finite-horizon}

Theorem~\ref{thm:reversible-reconstruction} reads the coordinates off the limiting boundaries.
At a finite horizon the observed boundaries are shifted, and the shift has an explicit form,
uniformly over a class of chains and starts. Throughout, the connected graph $E$ on $[d]$ and
the numbers $0<b\le R$ and $0<\mu_*\le1/d$ are fixed. Let $\mathcal C$ be the set of reversible
irreducible generators on $E$ with $\|S\|_2\le R$ and $\sqrt{w_{ij}}\ge b$ on every edge, and let
$\mathcal M$ be the set of starts with $\min_i\mu_i\ge\mu_*$. For a pair $(i,j)$ and $\xi\in\R$
let $v^{ij}(\xi)$ be the design field of Section~\ref{sec:reversible-reconstruction}, and let
$\mathcal V$ be the compact set of these fields with $|\xi|\le3R$.

The fixed-chain face maxima are Theorem~\ref{thm:facemax}. The finite horizon also needs
constants that do not depend on the chain and the start in the class. This is the last result
we import from~\cite{paperC1}. There it holds for an arbitrary compact set of fields, and we apply
it with $\mathcal V$. Our edge bound $b$ is the bound $w_*$ of~\cite{paperC1}, and our $S$ is its
$S_V$.

\begin{proposition}[uniform face maxima on the class, {\cite[Proposition~5.12]{paperC1}}]\label{prop:uniform-facemax}
Let $E$ induce a connected graph on the face $F$, with $k=|F|\ge2$. There are $N_0$ and $C$,
depending only on $E$, $R$, $b$ and $\mu_*$, such that for all $Q\in\mathcal C$, $\mu\in\mathcal M$
and $v\in\mathcal V$ the following hold.
\begin{enumerate}[(a)]
\item For all $N\ge2$ and $T>0$ with $TR\le N/2$,
$M_F(T,v)\le C(T/N)^{k-1}e^{-T\lambda_F(v)+CT^2/N}$.
\item For $N\ge N_0$ and $1\le T\le N^{1/3}$, we have $C^{-1}\le K_F(v)\le C$ and
\[
 \biggl|\frac{M_F(T,v)}{K_F(v)\,N^{-(k-1)}T^{(k-1)/2}e^{-T\lambda_F(v)}}-1\biggr|
 \le C\Bigl(\frac1T+\frac{T^2}N\Bigr).
\]
\end{enumerate}
For a singleton, $M_{\{i\}}(T,v)$ is given exactly by~\eqref{eq:vertexmax}.
\end{proposition}

The proof in~\cite{paperC1} is a compactness argument. On an edge $q_{ij}\ge b^2/R$, so the
class, the starts and the fields form a compact set on which every connected face stays
irreducible, and every constant in the fixed-chain proofs of Lemma~\ref{lem:upper} and
Theorem~\ref{thm:facemax} is continuous there.

\begin{theorem}[finite-horizon recovery on the class]\label{thm:finite-horizon}
Run the design of Theorem~\ref{thm:reversible-reconstruction} at horizon $N$, with $L=\log N$ and
$T=\tau L$, and let $\mathcal W_N(\tau,v)$ be the set of supports of the maximizers of
$n\mapsto P_N(n)e^{Tv\cdot n/N}$. For an ordered pair $(i,j)$ with $\{i,j\}\in E$ put
$v=v^{ij}(\bar\xi)$,
\[
 \tau^-_N=\sup\Bigl\{\tau\in[\tfrac1{4R},\tau_{\max}]:\mathcal W_N(\tau',v)=\{\{i\}\}\text{ for all }\tau'\in[\tfrac1{4R},\tau]\Bigr\},
\]
\[
 \tau^+_N=\inf\Bigl\{\tau\in[\tfrac1{4R},\tau_{\max}]:\mathcal W_N(\tau',v)=\{\{i,j\}\}\text{ for all }\tau'\in[\tau,\tau_{\max}]\Bigr\},
\]
let $t_{ij,N}$ be any number in $[\tau^-_N,\tau^+_N]$, and let $K_{ij}=K_{\{i,j\}}(v)$. Put
$\varepsilon_N=\log L/(2L)$. There are $N_0$ and $C$, depending only on $E$, $R$, $b$ and
$\mu_*$, such that for all $Q\in\mathcal C$, $\mu\in\mathcal M$ and $N\ge N_0$ the following
hold.
\begin{enumerate}[(a)]
\item (Singleton switches.) Let $i\ne1$. For every $\xi\in[-3R,3R]$ other than
\[
 x_{i,N}=\frac1{\tau_0L}\Bigl[\log\frac{\mu_1}{\mu_i}+(N-1)\log\frac{N-\tau_0Lq_1}{N-\tau_0Lq_i}\Bigr],
\]
$\mathcal W_N(\tau_0,v^{i1}(\xi))$ is $\{\{1\}\}$ for $\xi<x_{i,N}$ and $\{\{i\}\}$ for
$\xi>x_{i,N}$. Moreover $|x_{i,N}-x_i-2R\log(\mu_1/\mu_i)/L|\le CL/N$, so
$|x_{i,N}-x_i|\le2R\log(1/\mu_*)/L+CL/N$.
\item (Pair entries.) We have $\tau^+_N-\tau^-_N\le C/L^2$ and
\[
 \Bigl|t_{ij,N}-t_{ij}(1-\varepsilon_N)-\frac{c_{ij}}L\Bigr|\le C\,\frac{\log L}{L^2},\qquad
 c_{ij}=-t_{ij}\Bigl(\frac12\log t_{ij}+\log\frac{K_{ij}}{\mu_i}\Bigr),\qquad|c_{ij}|\le C .
\]
\item (Reconstruction.) Let $\widehat Q_N$ be the reconstruction of
Theorem~\ref{thm:reversible-reconstruction}(c) from the values $x_{i,N}$ and $t_{ij,N}$, and
$\widetilde Q_N$ the one from $x_{i,N}$ and $t_{ij,N}/(1-\varepsilon_N)$. Both are reversible
generators on $E$, and
\[
 \|\widehat Q_N-Q\|_2\le C\,\frac{\log L}L,\qquad\|\widetilde Q_N-Q\|_2\le\frac CL .
\]
\end{enumerate}
\end{theorem}

The threshold $N_0$ is large. The design runs up to $\tau=\tau_{\max}=8R/b^2$, and $I+(T/N)Q$
is stochastic on the whole range exactly when $N\ge \tau_{\max}q_{\max}\log N$. So the least such $N$ has
$Nb^2/\log N=8Rq_{\max}$ up to rounding, and it grows like $b^{-2}\log(1/b)$ when $q_{\max}$ stays
of order $R$. Section~\ref{sec:verify} computes it for 2 families. The proof compares the exact
singleton atoms with the uniform face maximum of the pair and solves the resulting crossing
equation for $T$.

\begin{proof}
All constants below depend only on $E$, $R$, $b$ and $\mu_*$. On the ranges of the design
$T\in[L/(4R),\tau_{\max}L]$, so for $N\ge N_1$ we have $TR\le N/2$, $T\le N^{1/3}$ and
$1/T+T^2/N\le5R/L$, and Proposition~\ref{prop:uniform-facemax} applies.

\emph{Suppression.} Fix a pair and a point of the ranges, and let $m^*$ be the least score among
the irreducible faces inside the pair. The proof of Lemma~\ref{lem:suppression} uses only $R$,
$b$ and $\bar v$, so each connected face $G'$ outside the pair has
$\Phi_{\tau,v}(G')\ge m^*+\frac34$, and Proposition~\ref{prop:uniform-facemax}(a) or the singleton
atom gives $M_{G'}\le CT^{|G'|-1}N^{-\Phi_{\tau,v}(G')}$. The best face inside the pair has face
maximum at least $C^{-1}N^{-m^*}$, by the singleton atom with $\mu_i\ge\mu_*$ or by
Proposition~\ref{prop:uniform-facemax}(b). Faces that are not connected carry no mass. So for
$N\ge N_0$ every maximizer has its support inside the pair.

\emph{(a)} At $\tau_0$ the pair $\{i,1\}$, if it is an edge, has score at least the better
singleton score plus $\frac12$, by the proof of Theorem~\ref{thm:reversible-reconstruction}(a), so
as above its face maximum is smaller than $\max(M_{\{1\}},M_{\{i\}})$ by a factor $CLN^{-1/2}$.
The two singleton atoms are exact, and
\[
 \log M_{\{i\}}-\log M_{\{1\}}=\log\frac{\mu_i}{\mu_1}+(N-1)\log\frac{N-\tau_0Lq_i}{N-\tau_0Lq_1}+\tau_0L\,\xi ,
\]
which is affine and increasing in $\xi$ with zero $x_{i,N}$. This gives the claim on
$\mathcal W_N$. Since $|\log(1-y)+y|\le y^2$ for $0\le y\le\frac12$, the middle term is
$-\tau_0L(q_i-q_1)+O(L/N)+O(R^2L^2/N)$, and dividing by $\tau_0L=L/(2R)$ gives the expansion. For
large $N$ it gives $x_{i,N}\in(-3R,3R)$, since $|x_i|\le R$.

\emph{(b)} Let $\{i,j\}$ be an edge and $v=v^{ij}(\bar\xi)$. The exact atoms give
$\log M_{\{j\}}-\log M_{\{i\}}\le\log(1/\mu_*)-TR+1$ with $TR\ge L/4$, so $\{j\}$ loses to
$\{i\}$. By the proof of Theorem~\ref{thm:reversible-reconstruction}(b),
$\lambda_{\{i\}}(v)-\lambda_{\{i,j\}}(v)=1/t_{ij}$, and the exact atom gives
$\log M_{\{i\}}=\log\mu_i-T(q_i-\bar\xi)+O(1/L)$. So Proposition~\ref{prop:uniform-facemax}(b)
with $k=2$ gives, on the whole range,
\[
 \Lambda_N(T)=\log M_{\{i,j\}}-\log M_{\{i\}},\qquad\bigl|\Lambda_N(T)+L-\Xi(T)-\kappa\bigr|\le\frac{C_5}L,
\]
with $\Xi(T)=T/t_{ij}+\frac12\log T$ and $\kappa=\log(K_{ij}/\mu_i)$. Here $|\kappa|\le C$, since
$K_{ij}$ is bounded above and below and $\mu_i\ge\mu_*$. The function $\Xi$ increases with
$\Xi'\ge1/t_{ij}\ge2/\tau_{\max}$. Let $T^\circ$ solve $\Xi(T^\circ)=L-\kappa$. Then $\Lambda_N(T)<0$ for
$T\le T^\circ-C_5\tau_{\max}/L$ and $\Lambda_N(T)>0$ for $T\ge T^\circ+C_5\tau_{\max}/L$. As $t_{ij}\in[1/R,\tau_{\max}/2)$, the
window $|T-T^\circ|<C_5\tau_{\max}/L$ lies inside the range for large $N$. With the suppression step, every
maximizer has support $\{i\}$ below the window and $\{i,j\}$ above it. Hence
$T^\circ/L-C_5\tau_{\max}/L^2\le\tau^-_N\le\tau^+_N\le T^\circ/L+C_5\tau_{\max}/L^2$, which is the first claim, and
$t_{ij,N}=T^\circ/L+O(1/L^2)$. To solve for $T^\circ$ put $\tau^\circ=T^\circ/L$. Then
$\tau^\circ=t_{ij}\bigl(1-(\frac12\log T^\circ+\kappa)/L\bigr)$, so
$\tau^\circ/t_{ij}=1+O(\log L/L)$ and $\log T^\circ=\log L+\log t_{ij}+O(\log L/L)$. Putting this
back gives $\tau^\circ=t_{ij}(1-\varepsilon_N)+c_{ij}/L+O(\log L/L^2)$, uniformly, since
$t_{ij}\le \tau_{\max}/2$ and $|\kappa|\le C$.

\emph{(c)} By (a) and (b) the dimensionless coordinates $x_{i,N}/R$ and $Rt_{ij,N}$ differ from
$x_i/R$ and $Rt_{ij}$ by at most $C\log L/L$, and after the correction by at most $C/L$, since
$t_{ij,N}/(1-\varepsilon_N)-t_{ij}=(c_{ij}/L+O(\log L/L^2))/(1-\varepsilon_N)$.
Theorem~\ref{thm:stability} gives its admissible error size $\varepsilon_0$ and its constant as
functions of $d$ and $b/R$ alone. So for $N\ge N_0$ these errors are below $\varepsilon_0$, and
that theorem gives both bounds and the generator property.
\end{proof}

\begin{remark}[a start that defeats every horizon]\label{rem:finite-start}
The bound $\min_i\mu_i\ge\mu_*$ cannot be dropped, since the start enters through the singleton
switches as $2R\log(\mu_1/\mu_i)/L$. Let $d=2$, $q_{12}=q_{21}=1$, $\tau=\frac14$ and $v=(0,1)$.
The scores are $\Phi(\{2\})=0$, $\Phi(\{1\})=\frac14$ and
$\Phi(\{1,2\})=1+\frac14\cdot\frac{1-\sqrt5}2>\frac34$, so for a fixed start $\{2\}$ wins for large
$N$. Both singleton atoms carry the factor $(1-T/N)^{N-1}$, so
$M_{\{1\}}/M_{\{2\}}=(\mu_1/\mu_2)e^{-T}=(\mu_1/\mu_2)N^{-1/4}$ exactly. For the start
$\mu_2=1/(1+2N^{1/4})$ this ratio is $2$ at every $N$. The law is linear in the start, so
Lemma~\ref{lem:upper} for the two starts at a single state gives
$M_{\{1,2\}}\le CLN^{-1/2}M_{\{1\}}$ with $C$ free of $\mu$. So $\{1\}$
wins at every large $N$, against the limiting prediction. The switch between $\{1\}$ and $\{2\}$ in
the variable $v_2$ sits at $\log(\mu_1/\mu_2)/(\tau L)=1+4\log2/L$ instead of at $q_2-q_1=0$. On
$\mathcal M$ the same shift is at most $\log(1/\mu_*)/(\tau L)$, uniformly. For the pair
entries the start enters through the term $\log(K_{ij}/\mu_i)$ of $c_{ij}$ in
Theorem~\ref{thm:finite-horizon}(b), which is unbounded as $\mu_i\to0$.
\end{remark}

\begin{remark}[the size of the correction]\label{rem:finite-size}
On the two-state class with $R=3$ and $b=1$, the chain $q_1=q_2=1$ with the uniform start has
$t_{12}=6.16$ and $c_{12}=-16.5$, and at $N=10^7$ the exact law gives $t_{12,N}=4.65$, against
$5.63$ for $t_{12}(1-\varepsilon_N)$ and $4.60$ for the prediction of (b). So the term
$c_{ij}/L$, which no correction fixed in advance removes, dominates $t_{ij}\varepsilon_N$ at
accessible $N$.
\end{remark}
